\documentclass[11pt,a4paper]{article}
\usepackage[margin=1.5cm,top=1.2cm,bottom=1.2cm]{geometry}
\usepackage[bidi=basic,english]{babel}
\babelprovide[import]{arabic}
\babelprovide[import]{french}
\babelfont{rm}{Latin Modern Roman}
\babelfont[arabic]{rm}[Scale=1.12]{Amiri}
\usepackage{xcolor}
\definecolor{navy}{RGB}{20,45,95}
\pagestyle{empty}
\newcommand{\en}[1]{\mbox{\textdir TLT\relax\foreignlanguage{english}{#1}}}
\setlength{\parindent}{0pt}
\newcommand{\sect}[1]{\par\vspace{9pt}{\color{navy}\bfseries\large #1}\par\vspace{1pt}{\color{navy}\hrule height 0.6pt}\vspace{4pt}}
\begin{document}
\begin{center}
{\small University of Manouba -- National School of Computer Sciences (ENSI) -- HANALab}\\[3pt]
{\large\bfseries\color{navy} Deep Learning-Based Prediction of Disease Progression Using Medical
Imaging: An Application to Diabetology}\\[3pt]
{\small Doctoral thesis in Computer Science (Artificial Intelligence) -- Nader \textsc{Belhadj}\\
Supervisor: Prof.\ Lassaad Latrach -- Co-supervisor: Prof.\ Ridha Ghayoula -- Co-advisor: Mohamed Amine Mezghich}
\end{center}

\sect{Summary}
Diabetic retinopathy (DR) is a leading cause of preventable blindness. This thesis addresses
three limitations of automated DR grading: no relational context between patients, no use
of the topology of the retinal vessel network, and no safety mechanism allowing the system
to abstain on ambiguous images. It proposes TOPORET, an inductive population graph whose
edges combine visual and topological (persistent-homology) similarity, and then DOTS, which
adds the RETFound foundation model, CORAL ordinal regression and entropy-based abstention.
Pre-specified tests support the contribution of the graph (H1: $t_4=3.82$, $p=0.019$) and
the association of topological features with severity (H2: $F=84.3$, $p<0.001$). On
$10{,}000$ held-out images, DOTS reaches QWK $=0.951$, a severe-DR non-referral rate of
$1.9\%$ (95\% CI $1.2$--$2.8\%$) and $93.8\%$ coverage; none of the six reduced systems
reaches the $2\%$ target. ``Disease progression'' is addressed as cross-sectional ordinal
grading. The evidence is retrospective: it is not a clinical validation, and prospective
multi-site studies are the next step. Supported by eleven publications.

\sect{Résumé}
\begin{otherlanguage}{french}
La rétinopathie diabétique (RD) est une cause majeure de cécité évitable. Cette thèse
traite trois limites du classement automatique de la RD : l'absence de contexte
relationnel entre patients, l'absence d'exploitation de la topologie du réseau vasculaire
rétinien, et l'absence d'un mécanisme de sécurité permettant de s'abstenir sur les images
ambiguës. Elle propose TOPORET, un graphe de population inductif dont les arêtes combinent
similarité visuelle et topologique (homologie persistante), puis DOTS, qui ajoute le modèle
de fondation RETFound, la régression ordinale CORAL et l'abstention par entropie. Des tests
fixés à l'avance soutiennent l'apport du graphe (H1 : $t_4=3{,}82$, $p=0{,}019$) et
l'association des descripteurs topologiques avec la sévérité (H2 : $F=84{,}3$,
$p<0{,}001$). Sur $10\,000$ images de test, DOTS atteint un QWK de $0{,}951$, un taux de
non-référence des cas sévères de $1{,}9\,\%$ (IC à 95\,\% : $1{,}2$--$2{,}8\,\%$) et une
couverture de $93{,}8\,\%$ ; aucun des six systèmes réduits n'atteint la cible de $2\,\%$.
La « progression de la maladie » est abordée par le grade ordinal, de façon transversale.
Les résultats sont rétrospectifs : ils ne constituent pas une validation clinique, et des
études prospectives multicentriques en sont la prochaine étape. Onze publications
soutiennent ce travail.
\end{otherlanguage}

\sect{\foreignlanguage{arabic}{ملخص}}
\begin{otherlanguage}{arabic}
يُعدّ اعتلال الشبكية السكري من أهم أسباب العمى الذي يمكن تفاديه. تعالج هذه الأطروحة ثلاث
نقائص في أنظمة التصنيف الآلي لاعتلال الشبكية السكري: غياب السياق العلائقي بين المرضى، وعدم
استغلال طوبولوجيا الشبكة الوعائية للشبكية، وغياب آلية أمان تسمح للنظام بالامتناع عن القرار في
الصور الغامضة. تقترح الأطروحة نظام
\en{TOPORET}،
وهو مخطط سكاني استقرائي تجمع روابطه بين التشابه البصري والتشابه الطوبولوجي (التماثل المستمر)، ثم
نظام
\en{DOTS}
الذي يضيف النموذج الأساسي
\en{RETFound}
والانحدار الترتيبي
\en{CORAL}
والامتناع القائم على الإنتروبيا. تدعم اختبارات محددة مسبقاً مساهمة المخطط
(الفرضية \en{H1}: \en{$t_4=3.82$}، \en{$p=0.019$})
وارتباط السمات الطوبولوجية بشدة المرض
(الفرضية \en{H2}: \en{$F=84.3$}، \en{$p<0.001$}).
على
\en{10,000}
صورة اختبار، يحقق
\en{DOTS}
معامل
\en{QWK = 0.951}
ونسبة عدم إحالة للحالات الشديدة قدرها
\en{1.9\%} (مجال الثقة \en{95\%}: من \en{1.2\%} إلى \en{2.8\%})
مع تغطية قدرها
\en{93.8\%}،
ولا يبلغ أيّ من الأنظمة المختزلة الستة الهدف المحدد
(\en{2\%}).
يُتناول «تطور المرض» من خلال التصنيف الترتيبي لدرجة الشدة بصورة مقطعية. النتائج استعادية ولا
تمثل تحققاً سريرياً، والخطوة التالية هي دراسات استباقية متعددة المراكز. يستند هذا العمل إلى
إحدى عشرة منشورة علمية.
\end{otherlanguage}

\vfill
{\footnotesize\color{gray}\centering Keywords / Mots-clés: diabetic retinopathy; graph neural networks; persistent homology;
ordinal regression; selective prediction; safety-oriented evaluation.\par}
\end{document}
